\section{Практичне використання криптографічних бібліотек}
\subsection{Симетричне автентифіковане шифрування}

Встановлення бібліотеки.
\begin{lstlisting}[language=Python]
import os
from cryptography.hazmat.primitives.ciphers.aead import AESGCM
from cryptography.exceptions import InvalidTag
\end{lstlisting}

\subsubsection{Генерація криптографічно стійкого ключа}

\begin{lstlisting}[language=Python]
key = AESGCM.generate_key(bit_length=256)
aesgcm = AESGCM(key)

message = b"Hello, my name is Mariia"
\end{lstlisting}

\subsubsection{Шифрування тестового повідомлення}

\begin{lstlisting}[language=Python]
nonce = os.urandom(12)
ciphertext = aesgcm.encrypt(nonce, message, None)

print("Message:", message)
print("Key:", key.hex())
print("Nonce:", nonce.hex())
print("Ciphertext:", ciphertext.hex())

\end{lstlisting}

\subsubsection{Розшифрування повідомлення}

\begin{lstlisting}[language=Python]
decrypted = aesgcm.decrypt(nonce, ciphertext, None)

print("\nDecrypted message:", decrypted)
\end{lstlisting}

\subsubsection{Перевірка цілісності після зміни шифротексту}

\begin{lstlisting}[language=Python]
modified_ciphertext = bytearray(ciphertext)
modified_ciphertext[0] ^= 1
modified_ciphertext = bytes(modified_ciphertext)

print("\nModified ciphertext:", modified_ciphertext.hex())

try:
    aesgcm.decrypt(nonce, modified_ciphertext, None)
    print("Decryption successful")
except InvalidTag:
    print("Error: authentication tag verification failed")
\end{lstlisting}

\subsubsection{Вивід}
У результаті виконання програми було згенеровано 256-бітний ключ, 12-байтовий nonce та отримано шифротекст тестового повідомлення.

\begin{lstlisting}[language=Python]
Message: b'Hello, my name is Mariia'
Key: fc5febdc3df5cd09ecb790ebea5651e7dd441783d988c1e643e58acc8db8f25d
Nonce: 9f72c8ca0d7e717102b1c728
Ciphertext: 1e256814693b724c81fb72b20cb1ee6bcae3db883a2a5fb36416e6af01e720ed626effe2d343afd2

Decrypted message: b'Hello, my name is Mariia'

Modified ciphertext: 1f256814693b724c81fb72b20cb1ee6bcae3db883a2a5fb36416e6af01e720ed626effe2d343afd2
Error: authentication tag verification failed
\end{lstlisting}

Після зміни одного байта шифротексту його розшифрування завершилося помилкою InvalidTag. Це свідчить про виявлення порушення цілісності зашифрованих даних.

\subsubsection{Призначення ключа, nonce та authentication tag}

У використаному алгоритмі AES-GCM ключ є секретним значенням, за допомогою якого виконується шифрування 
та розшифрування повідомлення. У цьому прикладі використовується ключ довжиною 256 біт.

Nonce --- це значення, яке використовується разом із ключем під час шифрування. 
У нашому випадку nonce має довжину 12 байт (Для AES-GCM рекомендована довжина nonce становить 12 байт (96 біт))
і генерується за допомогою криптографічно стійкого генератора випадкових чисел.
Для AES-GCM важливо, щоб nonce не повторювався при використанні одного й того самого ключа. 
Повторне використання nonce з тим самим ключем є небезпечним, оскільки може призвести до розкриття інформації про відкритий текст та порушення автентифікаційного захисту. 
Тому для кожного нового повідомлення необхідно використовувати новий nonce.

Authentication tag --- це тег, який генерується під час шифрування та використовується для 
перевірки цілісності й автентичності шифротексту. У використаній бібліотеці він формується автоматично.

Після зміни одного байта шифротексту перевірка authentication tag не проходить, тому виникає  \texttt{InvalidTag}. 
Це показує,що AES-GCM забезпечує не тільки конфіденційність даних, а й перевірку їх цілісності.

\subsection{Цифровий підпис}

\subsubsection{Генерація ключової пари}

% Код та результат

\begin{lstlisting}[language=Python]
# Bude cod 
\end{lstlisting}


\subsubsection{Створення цифрового підпису}

% Код та результат

\subsubsection{Перевірка цифрового підпису}

% Код та результат

\subsubsection{Перевірка підпису після зміни повідомлення}

% Код та результат
